\section{Software Tools}\label{sec:software_tools}

This section provides an overview of the software tools used in this research.

\subsection{Web Crawling and Data Collection}\label{subsec:web-crawling-and-data-collection}

\begin{itemize}
    \item \textbf{Playwright}\footnote{\url{https://playwright.dev/}} was used to run headless browsers to render web pages and save the HTML content for further processing.
    Playwright provides several useful features such as pre-navigation and post-navigation hooks, which allow for the execution of custom code before and after page navigation.
    This is particularly useful for handling dynamic content and ensuring that all necessary data is captured.
    Another useful feature of Playwright is its ability to handle multiple browser contexts, which allows for the simulation of different user sessions and the collection of data from multiple sources simultaneously.
    With this we can load multiple pages in parallel, which significantly speeds up the data collection process.
    Playwright also provides a powerful API for interacting with web pages, which allows for the automation of tasks such as hiding certain elements, clicking buttons, and filling out forms.
    Finally, Playwright also allows to define a target concurrency level as well as automatic throttling and retries, which helps to avoid overloading the target servers and ensures that the data collection process is efficient and reliable.

    \item \textbf{Amazon Simple Queue Service (SQS)}\footnote{\url{https://aws.amazon.com/sqs/}} was used to manage the queue of URLs to be crawled.
    SQS is a fully managed message queuing service that enables the decoupling and scaling of microservices, distributed systems, and serverless applications.
    SQS provides several useful features such as message visibility timeout, which allows for the temporary hiding of messages from other consumers while they are being processed, and message retention, which allows for the storage of messages for a specified period of time.
    Another useful feature was the ability to set up dead-letter queues, which allows for the handling of messages that cannot be processed successfully.
\end{itemize}

\subsection{Boilerplate Removal}\label{subsec:boilerplate_removal}

\begin{itemize}
    \item \textbf{BeautifulSoup}\footnote{\url{https://www.crummy.com/software/BeautifulSoup/}} is a Python library for parsing HTML and XML documents.
    This library was used as it offers a simple and intuitive API for navigating and manipulating the parse tree, which makes it easy to extract and modify the content of web pages.
    BeautifulSoup also provides several useful features for handling malformed or poorly structured HTML, which is common in web pages, and allows for the extraction of specific elements or attributes from the HTML content.
\end{itemize}

\subsection{Proof-Of-Concept Application}\label{subsec:proof-of-concept-application}

\begin{itemize}
    \item \textbf{Next.js}\footnote{\url{https://nextjs.org/}} was used to build the front-end of the proof-of-concept application.
    Next.js is a popular React framework that provides several useful features for building web applications, such as server-side rendering, static site generation, and automatic code splitting.
    These features help to improve the performance and scalability of the application, as well as providing a better developer experience.
\end{itemize}

\subsection{Data Storage and Management}\label{subsec:data_storage_and_management}
\begin{itemize}
    \item \textbf{Amazon Simple Storage Service (S3)}\footnote{\url{https://aws.amazon.com/s3/}} was used to store the crawled HTML content, minimized html content, model training data and the trained model artifacts.
    S3 is a fully managed object storage service that provides several useful features such as high durability, scalability, and availability.
    S3 also provides several storage classes, which allow for the optimization of storage costs based on the access patterns and retention requirements of the data.
    Another useful feature of S3 is its integration with other AWS services, which allows for the seamless transfer of data between different components of the application.
    \item \textbf{Aurora PostgreSQL on Amazon Relational Database Service (RDS)}\footnote{\url{https://aws.amazon.com/rds/aurora/postgresql/}} was used to store the metadata of the crawled web pages, such as the URL, title, and description.
    PostgreSQL also offers support for advanced data types including vectors, which is useful for storing and querying the embeddings generated by the product matching model.
    \item \textbf{Google BigQuery}\footnote{\url{https://cloud.google.com/bigquery}} was used as the data source for the observability dashboard, which allows for the monitoring of the different components of the application and provides insights into the performance and health of the system.
    BigQuery is a fully managed, serverless data warehouse that provides several useful features such as high scalability, fast query performance, and integration with other Google Cloud services.
    BigQuery was selected for this task, due to its deep integration with Google Data Studio, which allows for the easy creation of interactive dashboards and reports, performance with analytical workloads, and well as due to its generous free tier, which allows for the storage and querying of large datasets at almost no cost.
\end{itemize}

\subsection{Containerization and Orchestration}\label{subsec:containerization_and_orchestration}
\begin{itemize}
    \item \textbf{Docker}\footnote{\url{https://www.docker.com/}} was used to containerize the different components of the application, which allows for the isolation of dependencies and the easy deployment of the application across different environments.

    \item \textbf{Amazon Elastic Container Registry (ECR)}\footnote{\url{https://aws.amazon.com/ecr/}} was used to store the Docker images of the different components of the application, which allows for the easy distribution and deployment of the application.

    \item \textbf{Amazon Elastic Container Service (ECS)}\footnote{\url{https://aws.amazon.com/ecs/}} was used to orchestrate the deployment and scaling of the different components of the application, which allows for the efficient management of resources and the automatic scaling of the application based on demand.
    This was primarily used for the crawler deployment, which allows for the automatic scaling of the number of crawler instances based on the number of URLs in the SQS queue.
    Using spot instances for the crawler deployment also helps to reduce the cost of running the crawler, as spot instances are generally cheaper than on-demand instances.
\end{itemize}

\subsection{Model Training and Inference}\label{subsec:model_training_and_inference}
\begin{itemize}
    \item \textbf{Amazon SageMaker}\footnote{\url{https://aws.amazon.com/sagemaker/}} was used to train and deploy the product classification, information extraction, and product matching models.
    A unique feature offered by SageMaker is its support for running training jobs with hyperparameter tuning, which allows automatic optimization of model hyperparameters based on a specified objective metric.

    \item \textbf{Google Colab}\footnote{\url{https://colab.research.google.com/}} was used for prototyping and experimentation with different model architectures and training strategies.
    Google Colab provides a very user-friendly interface for running Jupyter notebooks in the cloud with access to free GPU and TPU resources, which allows for the efficient training of small machine learning models.

    \item \textbf{PyTorch}\footnote{\url{https://pytorch.org/}} was used to implement the machine learning models for product classification, information extraction, and product matching.
    PyTorch is a popular open-source machine learning framework that provides several useful features for building and training deep learning models, such as dynamic computation graphs, automatic differentiation, and a rich ecosystem of pre-trained models and libraries.
    PyTorch also provides several useful tools for model debugging and visualization, which helps to improve the interpretability and explainability of the models.

    \item \textbf{Hugging Face Transformers}\footnote{\url{https://huggingface.co/transformers/}} was used to implement the transformer-based models for product classification, information extraction, and product matching.
    Hugging Face Transformers is a popular open-source library that provides several pre-trained transformer models and tools for fine-tuning and deploying these models for various natural language processing tasks.
    The library also provides several useful features for model training and evaluation, such as data preprocessing, tokenization, and model evaluation metrics, which helps to improve the efficiency and effectiveness of the model training and evaluation process.
\end{itemize}

\subsection{Programming Languages}\label{subsec:programming_languages}

\begin{itemize}
    \item \textbf{TypeScript}\footnote{\url{https://www.typescriptlang.org/}} was used for several components of this project, such as the crawler and the web-interfaces.
    TypeScript is a common language used for web development (especially front-end), and it provides several useful features such as static typing, interfaces, and classes, which help to improve the maintainability and scalability of the codebase.
    TypeScript, as a superset of JavaScript, also allows for the use of modern JavaScript features and libraries, which helps to improve the performance and functionality of the application.

    \item \textbf{Python}\footnote{\url{https://www.python.org/}} was used for several components of this project, especially those related to machine learning and model training and inference.
    Python is a popular programming language for data science and machine learning, and it provides several useful libraries and frameworks for these tasks, such as NumPy, Pandas, Scikit-learn, TensorFlow, and PyTorch.
    These libraries and frameworks provide a wide range of functionality for data manipulation, analysis, and modeling, which helps to improve the efficiency and effectiveness of the machine learning components of the project.

    \item \textbf{Golang}\footnote{\url{https://golang.org/}} was used for several components of this project, especially those related to the backend APIs and data processing.
    Golang is a statically typed, compiled programming language that provides several useful features for building high-performance and scalable applications, such as concurrency, garbage collection, and a simple syntax.
    Golang also provides a rich standard library and a growing ecosystem of third-party libraries and frameworks, which helps to improve the efficiency and effectiveness of the backend components of the project.
    Golang compiled programs generally have a smaller memory footprint and faster execution time compared to interpreted languages like Python, which helps to improve the performance and scalability of the backend components of the project.
\end{itemize}

\subsection{Infrastructure as Code (Iac)}\label{subsec:infrastructure_as_code}
\begin{itemize}
    \item \textbf{Terraform}\footnote{\url{https://www.terraform.io/}} was used to manage the infrastructure of the application as code.
    Terraform is an open-source tool that allows for the definition and provisioning of infrastructure resources in a declarative manner, which helps to improve the reproducibility and maintainability of the infrastructure.
    Terraform also provides several useful features such as state management, which allows for the tracking of changes to the infrastructure over time, and modules, which allow for the reuse of infrastructure code across different projects and environments.
\end{itemize}

\subsection{Source Code Management}\label{subsec:source_code_management}
\begin{itemize}
    \item \textbf{GitHub}\footnote{\url{https://github.com}} was used to manage the source code for all the components of the application.
    GitHub provides several useful features such as version control, which allows for the tracking of changes to the source code over time, and collaboration, which allows for the sharing of code and the coordination of development efforts among multiple contributors.
    GitHub also provides several integrations with other tools and services, such as continuous integration and deployment (CI/CD) pipelines, which allows for the automatic testing and deployment of the application whenever changes are pushed to the repository.
\end{itemize}

\subsection{Continuous Integration and Continuous Deployment (Ci/cd)}\label{subsec:ci_cd}
\begin{itemize}
    \item \textbf{GitHub Actions}\footnote{\url{https://github.com/features/actions}} was used to automate the CI/CD pipelines for the application.
    GitHub Actions allows for the definition of workflows that can be triggered by various events, such as pushes to the repository or the creation of pull requests, which helps to improve the efficiency and reliability of the development process.
    This was used to orchestrate pipelines which automatically build, test, and deploy the different components of the application whenever changes are pushed to the repository.
\end{itemize}

\subsection{Other Tools}\label{subsec:other_tools}
\begin{itemize}
    \item \textbf{AWS Amplify}\footnote{\url{https://aws.amazon.com/amplify/}} was used to manage the deployment and hosting of the web applications related to this project.
    Amplify allows you to connect your application to GitHub repositories and automatically deploy the application whenever changes are pushed to the repository.
    \item \textbf{AWS Lambda}\footnote{\url{https://aws.amazon.com/lambda/}} was used to run serverless functions for various short-lived tasks such as handling API requests and managing the crawler queue.
    \item \textbf{AWS Step Functions}\footnote{\url{https://aws.amazon.com/step-functions/}} was used to orchestrate the different stages of the model inference pipeline, which allows for the efficient management of resources and automatic retries in case of failures.
    \item \textbf{Google Data Studio}\footnote{\url{https://datastudio.google.com/}} was used to create the observability dashboard for the application.
    Data Studio provides several useful features for creating interactive dashboards and reports, such as drag-and-drop interface, data connectors, and visualization options.
    \item \textbf{Google Datastream}\footnote{\url{https://cloud.google.com/datastream}} was used to replicate data from Aurora PostgreSQL to Google BigQuery, which allows efficient querying and analysis of data in the observability dashboard.
\end{itemize}